\begin{description}
\item[(a)] The Sun's total power can be calculated from the
  Stefan--Boltzmann law: \hfill(1 p)
  \begin{equation*}
  P = A\sigma T^4, \tag{7.1}
  \end{equation*}
  where $A$ is the emitting surface, $\sigma$ is the Stefan--Boltzmann
  constant and $T$ is the (effective) temperature.

  The total power emitted homogeneously and isotropically by the Sun
  ($P_\odot$) distributes uniformly on the surface area of a sphere with
  radius $d_\oplus = 1$\,au. Hence: \hfill(2 p)
  \begin{equation*}
  I_\odot = \frac{P_\odot}{4\pi d_\oplus^2}
          = \frac{\sigma 4\pi R_\odot^2 T_\odot^4}{4\pi d_\oplus^2}
          = \sigma\,\frac{R_\odot^2 T_\odot^4}{d_\oplus^2}, \tag{7.2}
  \end{equation*}
  where $R_\odot$ and $T_\odot$ denote the radius and effective temperature of
  the Sun, respectively.

  With the numerical values taken from the Table of constants:
  $I_\odot = 1365.9$\,W\,m$^{-2}$. \hfill(1 p)
\item[(b)] In the first days of January the Earth is in perihelion, while in
  the beginning of July it is in aphelion. Therefore, instead of the mean
  Earth--Sun distance (i.e.\ the semi-major axis of the Earth's orbit), now
  one should substitute the perihelion and aphelion distances, respectively.
  Therefore, one has to write that: \hfill(1 p)
  \begin{gather*}
  I_{\odot,\mathrm{perihelion}} = \frac{P_\odot}{4\pi(d_\oplus(1-e))^2}
    = \sigma \frac{R_\odot^2 T_\odot^4}{(d_\oplus(1-e))^2}
    = \frac{I_\odot}{(1-e)^2}, \tag{7.3}\\
  I_{\odot,\mathrm{aphelion}} = \frac{P_\odot}{4\pi(d_\oplus(1+e))^2}
    = \sigma \frac{R_\odot^2 T_\odot^4}{(d_\oplus(1+e))^2}
    = \frac{I_\odot}{(1+e)^2}, \tag{7.4}
  \end{gather*}
  where $e = 0.016\,710\,22$ is the eccentricity of the Earth's orbit (see the
  Table of constants).

  With numerical values:

  in early January: $I_{\odot,\mathrm{perihelion}} = 1412.7$\,W\,m$^{-2}$
  \hfill(1 p)

  in early July: $I_{\odot,\mathrm{aphelion}} = 1321.4$\,W\,m$^{-2}$
  \hfill(1 p)

  The ratio of the values:
  $I_{\odot,\mathrm{perihelion}}/I_{\odot,\mathrm{aphelion}} = 1.0691$. The
  perihelion irradiation is higher by about 7\%. \hfill(1 p)
\item[(c)] In this case the total power of the Sun becomes \hfill(2 p)
  \begin{equation*}
  P'_\odot = \sigma\left[T_\odot^4(A_\odot - A_{\mathrm{sp}})
    + T_{\mathrm{sp}}^4 A_{\mathrm{sp}}\right], \tag{7.5}
  \end{equation*}
  where $A_\odot$ stands for the solar surface, while $A_{\mathrm{sp}}$ is the
  area of the spot (as a circle).

  Thus, the solar constant takes the following form: \hfill(2 p)
  \begin{align*}
  I'_\odot &= \frac{P'_\odot}{4\pi d_\oplus^2}
    = \frac{\sigma\pi\left[T_\odot^4(4R_\odot^2 - R_{\mathrm{sp}}^2)
      + T_{\mathrm{sp}}^4 R_{\mathrm{sp}}^2\right]}{4\pi d_\oplus^2}
    = \frac{\sigma}{4 d_\oplus^2}\left[T_\odot^4(4R_\odot^2 - R_{\mathrm{sp}}^2)
      + T_{\mathrm{sp}}^4 R_{\mathrm{sp}}^2\right] = \tag{7.6}\\
  &= I_\odot\left[1 - \left(\frac{R_{\mathrm{sp}}}{2R_\odot}\right)^2
    + \left(\frac{R_{\mathrm{sp}}}{2R_\odot}\right)^2
      \left(\frac{T_{\mathrm{sp}}}{T_\odot}\right)^4\right]
   = I_\odot\left[1 - \left(\frac{R_{\mathrm{sp}}}{2R_\odot}\right)^2
      \left(1 - \left(\frac{T_{\mathrm{sp}}}{T_\odot}\right)^4\right)\right] =\\
  &= 0.9991\,I_\odot \tag*{(2 p)}
  \end{align*}
  With numerical value: $I'_\odot = 1364.6$\,W\,m$^{-2}$. \hfill(1 p)
\item[(d)] As we have seen, the ratio is
  \[
  A = \frac{I'_\odot}{I_\odot}
    = \frac{1364.64~\mathrm{W\,m^{-2}}}{1365.92~\mathrm{W\,m^{-2}}} = 0.9991,
  \]
  or, converting it into percentage, the solar constant in the case of the
  spotted Sun will be lower by 0.09\%.

  When this sunspot emerges in the surface of the Sun, solar emission no
  longer will be isotropic, i.e.\ uniform in every directions. When the
  sunspot is fully seen from the Earth, the ratio of the irradiation of the
  Earth in the spotted and unspotted cases will be equal to the ratio of the
  solar powers in the two cases, i.e.,
  \begin{gather*}
  P_{\odot,\mathrm{disk}} = \sigma T_\odot^4 A_\odot
    = \sigma\pi T_\odot^4 R_\odot^2 \tag{7.7}\\
  P'_{\odot,\mathrm{disk}} = \sigma\left[T_\odot^4(A_\odot - A_{\mathrm{sp}})
    + T_{\mathrm{sp}}^4 A_{\mathrm{sp}}\right]
    = \sigma\pi\left[T_\odot^4(R_\odot^2 - R_{\mathrm{sp}}^2)
    + T_{\mathrm{sp}}^4 R_{\mathrm{sp}}^2\right] \tag{7.8}
  \end{gather*}
  \hfill(1 p each)

  and therefore, the ratio of the irradiations (or the effective solar
  constants) can be written as: \hfill(2 p)
  \begin{equation*}
  A' = \frac{P'_{\odot,\mathrm{disk}}}{P_{\odot,\mathrm{disk}}}
     = \frac{T_\odot^4(R_\odot^2 - R_{\mathrm{sp}}^2)
       + T_{\mathrm{sp}}^4 R_{\mathrm{sp}}^2}{T_\odot^4 R_\odot^2}
     = 1 - \left(\frac{R_{\mathrm{sp}}}{R_\odot}\right)^2
       \left[1 - \left(\frac{T_{\mathrm{sp}}}{T_\odot}\right)^4\right] \tag{7.9}
  \end{equation*}
  With numerical values: $A' = 0.9963$. \hfill(1 p)

  Thus, when the spot directed towards the Earth the total irradiation of the
  Earth becomes weaker by 0.37\%.
  % sic: "when the spot directed towards" as in the original
\end{description}
